\clearpage
\VerdierPhysicalPageAnchor{0204}%
\begingroup\small\setlength{\parskip}{0pt}\setlength{\emergencystretch}{1.5em}
\setlength{\abovedisplayskip}{4pt plus 2pt minus 2pt}%
\setlength{\belowdisplayskip}{4pt plus 2pt minus 2pt}%

\VerdierSemanticAnchor{ch3d-2-4-4-b}
\noindent\textit{b) Désignons par $P$ l'ensemble des objets projectifs de
$\mathcal A$. La catégorie triangulée $K^{-}(P)$ est une sous-catégorie de
la catégorie des objets $Q$-libres à gauche. Le foncteur naturel :}
\[
K^{-}(P)\longrightarrow D(\mathcal A)
\]
\noindent\textit{est pleinement fidèle.}

\VerdierSemanticAnchor{ch3d-2-4-4-c}
\noindent\textit{c) Si tout objet de $\mathcal A$ est quotient d'objet
projectif (resp.\ et si de plus il existe un entier $n\in\mathbb N$ qui
majore la dimension projective des objets de $\mathcal A$), alors :}

\noindent\textit{1) Le foncteur naturel
$K^{-}(P)\longrightarrow D^{-}(\mathcal A)$ (resp.
$K^{*}(P)\longrightarrow D^{*}(\mathcal A)$) est une équivalence de
catégories. On désigne par
$L^{-}\colon D^{-}(\mathcal A)\longrightarrow K^{-}(P)$ (resp.
$L^{*}\colon D^{*}(\mathcal A)\longrightarrow K^{*}(P)$) un foncteur
quasi-inverse.}

\noindent\textit{2) Le foncteur
$L^{-}Q^{-}\colon K^{-}(\mathcal A)\longrightarrow K^{-}(P)$ (resp.
$L^{*}Q^{*}\colon K^{*}(\mathcal A)\longrightarrow K^{*}(P)$) est adjoint
à droite au foncteur d'inclusion :}
\[
K^{-}(P)\longrightarrow K^{-}(\mathcal A)
\qquad
(\text{resp. }K^{*}(P)\longrightarrow K^{*}(\mathcal A)).
\]
\noindent\textit{Le morphisme d'adjonction
$\operatorname{id}_{K^{-}(\mathcal A)}\longrightarrow L^{-}Q^{-}$
(resp.
$\operatorname{id}_{K^{*}(\mathcal A)}\longrightarrow L^{*}Q^{*}$) est
un ``quasi-isomorphisme''.}

\noindent\textit{3) Le foncteur $L^{-}$ (resp. $L^{*}$) est un adjoint à
gauche au foncteur $Q^{-}$ (resp. $Q^{*}$).}

\noindent\textit{4) Un objet de $K^{-}(\mathcal A)$ (resp.
$K^{*}(\mathcal A)$) est $Q$-libre à gauche si et seulement s'il est
isomorphe à un objet de $K^{-}(P)$ (resp. $K^{*}(P)$).}

\VerdierSemanticAnchor{ch3d-2-4-4-proof-a}
Démontrons a). Pour tout objet $Y$ de $\mathcal A$, on a :
\[
\operatorname{Hom}_{K(\mathcal A)}(X,Y[n])=0,
\qquad\text{pour }n\neq0,
\]
et par suite, comme $X$ est $Q$-libre à gauche (chap. II, 2.3.3), on a :
\[
\operatorname{Hom}_{D(\mathcal A)}(X,Y[n])=0,
\qquad\text{pour }n\neq0.
\]
Par ailleurs, on a
$\operatorname{Hom}_{D(\mathcal A)}(X,Y)
=\operatorname{Hom}_{\mathcal A}(X,Y)$ (1.2.9). Soit alors :

\VerdierSemanticAnchor{ch3d-2-4-4-proof-a-short-exact}
\[
S=0\longrightarrow Y'\xrightarrow{u}Y\xrightarrow{v}Y''
\longrightarrow0
\]
une suite exacte de $\mathcal A$. Le foncteur cohomologique canonique
$\operatorname{comp}(\mathcal A)\longrightarrow D(\mathcal A)$ (1.3.2)
fournit un triangle distingué
$(Y',Y,Y'',[u],[v],\delta(S))$. D'où, une suite exacte (chap. II, 1.2.1) :

\VerdierSemanticAnchor{ch3d-2-4-4-proof-a-hom-exact}
\[
\cdots\longrightarrow0\longrightarrow
\operatorname{Hom}_{D(\mathcal A)}(X,Y')\longrightarrow
\operatorname{Hom}_{D(\mathcal A)}(X,Y)\longrightarrow
\operatorname{Hom}_{D(\mathcal A)}(X,Y'')\longrightarrow0
\longrightarrow\cdots,
\]
i.e.\ une suite exacte :

\VerdierSemanticAnchor{ch3d-2-4-4-proof-a-hom-a-exact}
\[
0\longrightarrow\operatorname{Hom}_{\mathcal A}(X,Y')
\longrightarrow\operatorname{Hom}_{\mathcal A}(X,Y)
\longrightarrow\operatorname{Hom}_{\mathcal A}(X,Y'')
\longrightarrow0,
\]
\endgroup
